\subsection{Labeled Timed Petri Nets}
\label{sec:LTPN}

In this paper, we use Labeled Timed Petri Nets (LTPN) a subclass of Petri nets to represent the 
stochastic scheduling problem. 
A Petri net $\mathcal{N} = (E, P, F)$ is a directed graph consisting of a set of nodes, called places ($P$), and a set of transitions ($E$)~\cite{murata1989petri}. These nodes are connected by directed arcs, represented as the flow relation $F \subseteq (E \times P) \cup (P \times E)$. A Labeled Timed Petri Net (LTPN) is an extension, denoted as $\mathcal{N}^{LPTN} = (E, P, F, A, l, \tau)$, where $E$ contains two types of transitions: $E' \subseteq E$ as a set of timed transitions, and $E \backslash E'$ as a set of immediate transitions. The function $l : E' \rightarrow A$ is a labeling function that assigns labels from a set $A$ to the timed transitions, and $\tau : E' \rightarrow \mathbb{R}$ is a timing function that assigns durations to the timed transitions in $E'$.

\textcolor{blue}{Arik continues here.}
Finally, $P_r \subset P$ denotes the places that represent the resources involved in the process.

Figure~\ref{fig:TPN} is an LTPN such that the activities \act{Blood Draw}, \act{Visits}, \act{Examination}, \act{Chemo. Infusion} define the timed transitions $E'$ while \{\act{Inv$_1$},....\act{Inv$_n$}\} $ \subset E'\backslash E$ the immediate transitions.
The places are represented by circles, with those labeled \act{N}, \act{P}, and \act{IN} corresponding to $P_{r}$ the resources involved in the process: nurses (N), physicians (P), and infusion nurses (IN).
The places $p_{Start}$ and $p_{End}$ correspond to the start and end points of the project, respectively.

Given a transition $e \in E$, $\bullet e$ denotes the set of input places for a transition $e \in E$, which are the places $p \in P$ that have a directed arc from $p$ to $e$, i.e., $(p, e) \in F$. Similarly, $e \bullet$ indicates the set of output places, which are the places $p$ that have a directed arc from $t$ to $p$. Referring to the example in Figure~\ref{fig:TPN}, $\bullet$\act{Vitals} $=$ {$p5$} and \act{Vitals}$\bullet$ $=$ {$p7$, \act{N}}. At any time, a place can contain zero or more tokens, represented by black or colored dots, with the latter specifically representing resources. For instance, if \act{N} contains two dots, it means that two nurses are involved in the process.

The state of an LPTN is determined by the distribution of tokens over places and is referred to as its marking. A marking of an LPTN is a function $M : P \rightarrow \mathbb{N}$ that maps each place to the number of tokens it contains, and the states change discretely by adding or removing tokens from places. A transition $t$ is enabled at a marking $m$ iff each input place contains at least one token, i.e., $\forall p \in \bullet e, m(p) \ge 0$. An enabled immediate transition can fire, while a timed transition $e' \in E'$ depends on its duration $\tau(e')$ i.e., firing requires a certain amount of time.
Firing a transition $e$ at marking $m$ yields a new marking $m'$, denoted as $m \xrightarrow[]{e} m'$.
$m_0$ and $m_t$ indicate respectively the initial and the target marking of LTPN, then a sequence of transition firing $s = \langle e_1, .... e_n \rangle \in E$ that leads from $m_0$ to $m_n$, such that $m_0 \xrightarrow[]{e_1} ... \xrightarrow[]{e_n} m_n$. 

Figure~\ref{fig:TPN} represents the initial marking $m_0$, where only \act{Inv1} is enabled. Once that transition fires, the two tokens (the project token and the nurse token) merge into a single token, enabling the timed labelled transition \act{Blood Draw}. After the time $\tau($\act{Blood Draw}$)$ required to execute the transition, one token is produced for the \act{N} place (simulating the release of the resource) and one for another place.
The process continues with the parallel execution of the \act{Visits} and \act{Examination} activities, performed by a nurse and a physician, respectively. Finally, the physician decides whether to proceed with the \act{Chemo. Infusion} based on the previous activities performed. Hence, the LTPN in Figure~\ref{fig:TPN} represents two possible paths of patients in the process. 


%%% VERSION OF MOTIVATING EXAMPLE AND DEFINITION OF LTPN TOGETHER 
\begin{comment}
A Petri net $\mathcal{N} = (E, P, F)$ is a directed graph consisting of a set of nodes, called places ($P$), and a set of transitions ($E$)~\cite{murata1989petri}. These nodes are connected by directed arcs, represented as the flow relation $F \subseteq (E \times P) \cup (P \times E)$. A Labelled Timed Petri Net (LTPN) is an extension, denoted as $\mathcal{N}^{LPTN} = (E, P, F, A, l, \tau)$, where $E$ contains two types of transitions: $E' \subseteq E$ as a set of timed transitions, and $E \backslash E'$ as a set of immediate transitions. The function $l : E' \rightarrow A$ is a labelling function that assigns labels from a set $A$ to the timed transitions, and $\tau : E' \rightarrow \mathbb{R}$ is a timing function that assigns durations to the timed transitions in $E'$.

Consider the LPTN in Figure~\ref{fig:TPN}, which represents a possible path in a hospital process. 
It contains four labelled and timed transitions $E'$ (\act{Blood Draw}, \act{Visits}, \act{Examination}, \act{Chemo. Infusion}) represented by white rectangles, and seven immediate transitions represented by smaller black rectangles. The places are represented by circles, with those labelled \act{N}, \act{P}, and \act{IN} corresponding to the resources involved in the process: nurses (N), physicians (P), and infusion nurses (IN).
The places \act{Start} and \act{End} correspond to the start and end points of the project, respectively.

Given a transition $e \in E$, $\bullet e$ denotes the set of input places for a transition $e \in E$, which are the places $p \in P$ that have a directed arc from $p$ to $e$, i.e., $(p, e) \in F$. Similarly, $e \bullet$ indicates the set of output places, which are the places $p$ that have a directed arc from $t$ to $p$. Referring to the example in Figure~\ref{fig:TPN}, $\bullet$\act{Vitals} $=$ {$p5$} and \act{Vitals}$\bullet$ $=$ {$p7$, \act{N}}. At any time, a place can contain zero or more tokens, represented by black or colored dots, with the latter specifically representing resources. For instance, if \act{N} contains two dots, it means that two nurses are involved in the process.

The state of an LPTN is determined by the distribution of tokens over places and is referred to as its marking. A marking of an LPTN is a function $M : P \rightarrow \mathbb{N}$ that maps each place to the number of tokens it contains, and the states change discretely by adding or removing tokens from places. A transition $t$ is enabled at a marking $m$ iff each input place contains at least one token, i.e., $\forall p \in \bullet e, m(p) \ge 0$. An enabled immediate transition can fire, while a timed transition $e' \in E'$ depends on its duration $\tau(e')$ i.e., firing requires a certain amount of time.
Firing a transition $e$ at marking $m$ yields a new marking $m'$, denoted as $m \xrightarrow[]{e} m'$.
$m_0$ and $m_t$ indicate respectively the initial and the target marking of LTPN, then a sequence of transition firing $s = \langle e_1, .... e_n \rangle \in E$ that leads from $m_0$ to $m_n$, such that $m_0 \xrightarrow[]{e_1} ... \xrightarrow[]{e_n} m_n$. 

Figure~\ref{fig:TPN} represents the initial marking $m_0$, where only \act{Inv1} is enabled. Once that transition fires, the two tokens (the project token and the nurse token) merge into a single token, enabling the timed labelled transition \act{Blood Draw}. After the time $\tau($\act{Blood Draw}$)$ required to execute the transition, one token is produced for the \act{N} place (simulating the release of the resource) and one for another place.
The process continues with the parallel execution of the \act{Visits} and \act{Examination} activities, performed by a nurse and a physician, respectively. Finally, the physician decides whether to proceed with the \act{Chemo. Infusion} based on the previous activities performed. Hence, the LTPN in Figure~\ref{fig:TPN} represents two possible paths of patients in the process. 
\end{comment}

%A LTPNs is a tupla $\mathcal{N} = \langle E, E', P, F, A, l, \tau \rangle$ where $E$ is a finite set of transitions with $E'\subseteq E$ a set of timed transitions and $E \backslash E'$ a set of immediate transitions, $P$ is a finite set of places and $F \subseteq E \, x \,  P \cup P \, x \, E$ is a set of directed arcs, called the flow relation. $A$ is a set of activity labels, and $l \in E' \rightarrow A$ is a labeling function. Finally, $\tau$ is a function $E' \rightarrow \mathbb{R}$ that assigns durations to timed transitions.



\subsection{Mapping LPTN to S-RCMSPS}

Starting from the LPTN, we can directly derive the corresponding stochastic(S) RCMSPS. To achieve this, we employ the method proposed by~\cite{senderovich2019learning} and extend it to represent both stochastic durations and resource calendars.

Following the definition of the S-RCMSPS problem, we present the translation from an LPTN to the corresponding S-RCMSPS.

\begin{define}[LTPN to S-RCMSPS]
Given an LTPN and an initial marking, $(\mathcal{N}^{LPTN} = (E, P, F, A, l, \tau), m_0)$, the S-RCMSPS is created as follows:
\begin{itemize}
    \item A project $i \in \mathcal{I}$ is defined for each tokens in $p_{start}$, with $|p_{start}| = |\mathcal{I}|$
    \item $N_i$ of project $i$ is based on the path in $\mathcal{N}^{LPTN}$ that token $t$ has to follow to reach $p_{end}$. Specifically, from the sequence of firings of a token $t$, $s_t = \langle e_1, .... e_n \rangle \in E$ we select only the timed transitions $N_i = \{l(e') \; | \; \forall e' \in s_t\}$.
    \item The precedence relations inside $N_i$ is given by, $\Pi_i = \{ (n_{iw}, n_{iz}) \, | \, n_{iw} = l(e'_w) \, \wedge \, n_{iz} = l(e'_z) \, \wedge \, e'_w \rightarrow e'_z \}$ with $e'_w \rightarrow e'_z$ being true if $e'_w$ occurs before $e'_z$ in $\mathcal{N}^{LPTN}$
    \item The resource set is given by the set of resource place, $\mathcal{R} = P_r$ and $C_k$ is defined based on the number of tokens $|P_{r_k}|$ in the initial marking $m_0$
    \item Given $e'$ representing the activity $n_{ij}$, the consumption of resource $c_{ij}$ is defined by the input resource places, $\bullet e' \subseteq P_r$ and $d_{ij}$ correspond to $\tau(e')$
\end{itemize}
\end{define}



\begin{table*}[h]
\centering
\caption{Kendall Tau metrics: $\frac{(C-D)}{(C+D)}$ where $C$ represents the number of concordant pairs and $D$ represents the number of discordant pairs. $\tau_{0.1} = (7-4)/(7+4) = 0.27$}
\scalebox{0.90}{
\begin{tabular}{l@{\hskip 0.1in}c@{\hskip 0.1in}c@{\hskip 0.1in}c@{\hskip 0.1in}c@{\hskip 0.1in}c@{\hskip 0.1in}c@{\hskip 0.1in}c@{\hskip 0.1in}c@{\hskip 0.1in}c@{\hskip 0.1in}c@{\hskip 0.1in}}
\toprule
\textbf{Name} & \textbf{Size} & \textbf{Unc.} & \textbf{\#Sim} & $q_{value}$ & $D_{DET}$ & $D_{alpha}$ &
$\mu_{SIM}$ & $\sigma_{SIM}$ & Ranking$_{DET}$ & Ranking$_{PROB}$\\
& & \textbf{Level} &  & &  & & &  & & \\
\midrule

\multirow{10}{*}{abz6} & \multirow{10}{*}{10 X 10} & \multirow{10}{*}{0.1} & \multirow{10}{*}{1000} & \multirow{10}{*}{0.663}& 1099 & 1100.789 & 1075.088 & 16.006 & 11 &  11\\
&  &  &  &  & 1098 & 1032.129 & 1008.824 & 13.481 & 10 & 10\\
&  &  &  &  & 1020 & 1023.122 & 999.844 & 13.885 & 9 & 9\\
&  &  &  &  & 1009 & 1012.848 & 986.749 & 15.790 & 8 & 8\\
&  &  &  &  & 993 & 992.600 & 968.395 & 13.564 & 7 & \underline{6} \\
&  &  &  &  & 992 & 991.965 & 969.105 & 14.459 & 6 & \underline{7}\\
&  &  &  &  & 985 & 984.544 & 961.937 & 12.945 & 5 & 5\\
&  &  &  &  & 976 & 981.818 & 960.354 & 12.461 & 4 & 4\\
&  &  &  &  & 973 & 978.818 & 956.565 & 12.761 & 3 & \underline{2}\\
&  &  &  &  & 972 & 978.448 & 958.575 & 11.809 & 2 & \underline{3}\\
&  &  &  &  & \colorbox{yellow}{971} & \colorbox{yellow}{976.627} & 955.858 & 12.658 & 1 & 1\\
\midrule
\multirow{17}{*}{abz6} & \multirow{17}{*}{10 X 10} & \multirow{17}{*}{0.5} & \multirow{17}{*}{1000} & \multirow{17}{*}{0.663}& 1180 & 1209.659 & 1101.562 & 63.845 & 17 & \underline{15}\\
&  &  &  &  & 1166 & 1203.719 & 1095.272 & 64.500 & 16 & \underline{13}\\
&  &  &  &  & 1150 & 1194.403 & 1101.857 & 55.855 & 15 & \underline{11}\\
&  &  &  &  & 1146 & 1201.416 & 1102.451 & 57.850 & 14 & \underline{12}\\
&  &  &  &  & 1145 & 1211.513 & 1112.775 & 58.055 & 13 &  \underline{16}\\
&  &  &  &  & 1139 & 1211.560 & 1117.581 & 54.006 & 12 & \underline{17}\\
&  &  &  &  & 1138 & 1205.774 & 1114.074 & 54.002 & 11 & \underline{13}\\
&  &  &  &  & 1133 & 1179.867 & 1083.465 & 57.152 & 10 & 10\\
&  &  &  &  & 1131 & 1174.709 & 1077.325 & 55.202 & 9 & 9\\
&  &  &  &  & 1130 & 1172.610 & 1072.796 & 57.174 & 8 & \underline{5}\\
&  &  &  &  & 1125 & 1173.187 & 1079.894 & 53.504 & 7 & \underline{6}\\
&  &  &  &  & 1123 & 1165.490 & 1078.271 & 52.660 & 6 & \underline{2}\\
&  &  &  &  & 1121 & 1173.881 & 1079.883 & 56.342 & 5 & \underline{7}\\
&  &  &  &  & 1117 & 1169.416 & 1071.345 & 58.627 & 4 & \underline{3}\\
&  &  &  &  & 1116 & \colorbox{yellow}{1165.397} & 1062.345 & 61.568 & 3 & \underline{1}\\
&  &  &  &  & 1115 & 1170.846 & 1082.280 & 54.992 & 2 & \underline{4}\\
&  &  &  &  & \colorbox{yellow}{1114} & 1177.483 & 1079.376 & 55.982 & 1 & \underline{8}\\
\midrule
\multirow{12}{*}{abz6} & \multirow{12}{*}{10 X 10} & \multirow{12}{*}{1} & \multirow{12}{*}{1000} & \multirow{12}{*}{0.663}& 1328 & 1386.059 & 1227.822 & 88.073 & 12 & \underline{10}\\
&  &  &  &  & 1300 & 1352.320 & 1216.442 & 84.863 & 11 & \underline{2}\\
&  &  &  &  & 1282 & 1390.128 & 1239.988 & 85.501 & 10 & \underline{11}\\
&  &  &  &  & 1276 & \colorbox{yellow}{1346.128} & 1210.492 & 82.739 & 9 & \underline{1}\\
&  &  &  &  & 1268 & 1417.516 & 1263.462 & 85.722 & 8 & \underline{12}\\
&  &  &  &  & 1261 & 1376.925 & 1219.810 & 93.901 & 7 & \underline{7}\\
&  &  &  &  & 1256 & 1384.173 & 1219.903 & 94.138 & 6 & \underline{9}\\
&  &  &  &  & 1244 & 1360.914 & 1207.348 & 86.688 & 5 & \underline{3}\\
&  &  &  &  & 1243 & 1376.867 & 1228.142 & 86.291 & 4 & \underline{8}\\
&  &  &  &  & 1242 & 1374.822 & 1228.975 & 86.330 & 3 & \underline{5}\\
&  &  &  &  & 1241 & 1374.132 & 1221.080 & 86.971 & 2 & \underline{4}\\
&  &  &  &  & \colorbox{yellow}{1240} & 1375.860 & 1222.953 & 89.360 & 1 & \underline{6}\\
\bottomrule
\end{tabular}}
%\caption{Cycle Time Distribution Distance and Conformance score}
\label{tab:resultNotCalendar}
\end{table*}